%%%%%%%%%%%%%%%%%%%%%%%%%%%%%%%%%%%%%%%
% Cover Letter - Магнит OMNI, Руководитель по развитию стратегии и проектного офиса
%%%%%%%%%%%%%%%%%%%%%%%%%%%%%%%%%%%%%%%

\documentclass[]{cover}
\usepackage{fancyhdr}

\pagestyle{fancy}
\fancyhf{}

\rfoot{Page \thepage \hspace{0pt}}
\thispagestyle{empty}
\renewcommand{\headrulewidth}{0pt}

\begin{document}

%%%%%%%%%%%%%%%%%%%%%%%%%%%%%%%%%%%%%%
%     TITLE NAME
%%%%%%%%%%%%%%%%%%%%%%%%%%%%%%%%%%%%%%
\namesection{}{\Huge{Алексей Орлов}}{  \href{mailto:lex.orlov78@yandex.ru}{lex.orlov78@yandex.ru} | +7~(901)~797-54-88 |  \urlstyle{same}\href{https://www.linkedin.com/in/alex-orlov-v/}{LinkedIn}
}

%%%%%%%%%%%%%%%%%%%%%%%%%%%%%%%%%%%%%%
%     MAIN COVER LETTER CONTENT
%%%%%%%%%%%%%%%%%%%%%%%%%%%%%%%%%%%%%%

\currentdate{\rutoday}
\lettercontent{Здравствуйте, команда Магнит OMNI,}

\lettercontent{Меня заинтересовала позиция Руководителя по развитию стратегии и проектного офиса. За последние годы я системно решал именно эту задачу - переводить цели руководства в измеримые инициативы и держать портфель проектов под контролем - сначала в Банке ВТБ, затем в ООО «ПочтаТех» (ИТ-контур логистики и доставки группы «Почта России»), что напрямую перекликается с доменом Магнит Доставка.}

\lettercontent{Магнит OMNI объединяет доставку, маркетплейс, программу лояльности и приложение в единую омниканальную экосистему - на стыке нескольких бизнес-направлений, где стратегия должна декомпозироваться в измеримый план, а не оставаться декларацией. Именно такие задачи меня интересуют.}

\lettercontent{За 6 лет я построил процесс стратегического квартального планирования и приоритизации инициатив в ООО «ПочтаТех», а в Банке ВТБ координировал портфельное планирование (PI Planning) на позиции ведущего коуча:}

{\raggedright\fontspec{DejaVu Sans}\fontsize{11pt}{13pt}\selectfont
\begin{itemize}
    \item \textbf{Декомпозиция целей и квартальное планирование:} спроектировал модель приоритизации инициатив и квартального планирования в ПочтаТех, связывающую цели руководства с измеримыми roadmap-планами, сократившую Time-to-Market примерно на 20\%
    \item \textbf{Портфель проектов и метрики:} настроил систему метрик (TTM, Cycle Time, SLA, план-факт по срокам и стоимости) для регулярного review портфеля перед CTO и руководителями направлений
    \item \textbf{Кросс-функциональная координация:} вёл кросс-функциональные ИТ-проекты в ритейле (М.Видео-Эльдорадо, Спортмастер) и координировал стейкхолдеров в ВТБ
    \item \textbf{ИИ в самом управлении проектом:} автоматизирую аналитику и отчётность Python-скриптами через REST API, использую Claude Code для подготовки материалов для руководства
\end{itemize}\par}
\vspace{6pt}

\lettercontent{У меня нет формального опыта стратегического консалтинга (Big-3/Big-4), но задачу, которую он обычно решает - разложить стратегию в измеримый план и держать портфель под контролем - я решал внутри компаний, с прямой ответственностью за результат. Техническое образование (к.т.н.) помогает предметно разговаривать с разработчиками, а опыт фасилитации и коучинга (ICF) - вести команды через изменения без потери доверия. Буду рад обсудить, как мой опыт может усилить проектный офис Магнит OMNI.}

\enlargethispage{2\baselineskip}
\begin{flushright}
\closing{С уважением,}

\signature{Алексей Орлов}
\end{flushright}
\end{document}
